\section{Displayed data II --- Mazur-style numerals for
$X_0(13)$}\label{sec:mazur}

The file
\texttt{Mazur\_X0\_13\_No\_Isogeny\_final.lean}
(\texttt{2f4e16e}; incremental
\href{https://github.com/DavidFox998/beal-level-26-foundations/actions/runs/35488699511}{35488699511},
$252$s) displays numerals associated with $X_0(13)$~\cite{mazur1978}.
None of them excludes a rational $13$-isogeny; see the honesty note
below.

\begin{displaythm}[Mazur numerals]
\label{disp:mazur}
\begin{itemize}[leftmargin=1.6em]
\item $\lvert\SL_2(\F_{13})\rvert=13\cdot 12\cdot 14=2184$
  (\texttt{SL2\_F13\_card});
\item $48<2184$ (\texttt{forty\_eight\_lt\_2184});
\item $13=2^2+3^2$ (\texttt{thirteen\_eq\_two\_squares});
\item $288/48=6$ (\texttt{two\_eighty\_eight\_div\_forty\_eight});
\item $(\{0,1\}:\mathrm{Finset}\,\mathbb{N}).\mathrm{card}=2$
  (\texttt{displayed\_cusps\_card}).
\end{itemize}
The first four identities are axiom-free \texttt{decide} theorems.
The cusp-card theorem uses $[\mathtt{propext},\mathtt{Quot.sound}]$.
\end{displaythm}

The single named theorem \texttt{\seqsplit{displayed\_X0\_13\_numerical\_data}}
re-exports the inhabited conjunction of Display~\ref{disp:mazur}
with the genus Nat $0$; despite the name, it does not state or
prove that a rational $13$-isogeny is excluded.

\begin{honesty}[Genus $0$ versus two cusps, and no isogeny exclusion]
$X_0(13)$ has genus $0$, so $X_0(13)(\Q)$ is infinite over $\Q$.
The pasted equality $\{2\text{ cusps}\}=X_0(13)(\Q)$ is
literature-false and is \emph{not} a theorem. The honest Lean
display is finite card~$2$ versus the genus-$0$ note
\texttt{X0\_13\_genus\_nat $\neq$ X0\_13\_cusp\_count}, packaged
in \texttt{\seqsplit{displayed\_X0\_13\_numerical\_data}}; it is not
a Mathlib proof that the rational points are infinite, and it is
not Mazur's isogeny theorem. In particular, the introduction's
classical strategy paragraph names Mazur's theorem for context only:
no Lean theorem in this repository excludes a rational $13$-isogeny
of the Frey curve, and the genus-$0$/infinitude fact recorded here
is logically unrelated to such an exclusion.

\emph{v31 rename.} Earlier releases (through \texttt{v0.30})
exported this same conjunction four further times, under the names
\texttt{X0\_13\_Q\_infinite},
\texttt{\seqsplit{frey\_no\_rational\_13\_isogeny}},
\texttt{\seqsplit{Serre\_non\_Borel\_mod13}}, and
\texttt{\seqsplit{mazur\_no\_Frey\_13\_isogeny}}, plus a fifth,
slightly smaller packaged theorem once named
\texttt{\seqsplit{Mazur\_X0\_13\_No\_Isogeny\_final}}. Those names invited a
reader to mistake the numeral display for an isogeny exclusion, a
non-Borel image statement, or Mazur's theorem itself. They are
removed as of this revision; there is no compatibility alias. The
differently-scoped parent names on
\texttt{Mazur\_X0\_13\_No\_Isogeny.lean} are untouched and remain
\texttt{def Prop}; it is that parent Prop, not any theorem in the
final-display file, that would correspond to an actual exclusion
statement.
\end{honesty}

Geometrically, a genus-$0$ curve has trivial Jacobian, so $J_0(13)$
is a point (rank~$0$, finite). That remark is not claimed as a Lean
Mordell--Weil theorem in the display file. Mazur's theorem that the
only rational $13$-isogenies of elliptic curves over $\Q$ arise from
the cusps remains \texttt{def Prop} on the parent, uninhabited.

Axioms: numerals $[\mathtt{propext}]$ or axiom-free;
package $[\mathtt{propext},\mathtt{Quot.sound}]$;
LLL no-go persistence
$[\mathtt{propext},\mathtt{Classical.choice},\mathtt{Quot.sound}]$.
No \texttt{native\_decide}, no \texttt{sorryAx}, no new axiom.
